\section{Extended Related Work} \label{app: rl}
This work both formalizes the challenges when applying contrastive reinforcement learning to manipulation and proposes a sampling-based strategy for leveraging interactions to mitigate those challenges. This places it in the intersection of work investigating CRL and robot learning for manipulation using interactions. 

\textbf{Contrastive Reinforcement Learning}: CRL is the application of contrastive learning~\citep{oord2018representation} to goal-conditioned reinforcement learning (GCRL), both as a classification~\citep{eysenbach2021clearning} and regression objective ~\citep{eysenbach2022contrastive}.  Recent work has expanded this towards metric learning~\citep{zheng2025multistep,myers2026offline}, planning as time series modeling~\citep{eysenbach2024inference}, language alignment~\citep{myers2023goal,myers2026temporal} and combinatorial reasoning and search~\citep{wang2026temporal,ziarko2026contrastive}. \methodname{} bears the closest similarity to~\citet{ziarko2026contrastive}, which also employs a resampling procedure, but \methodshort{} leverages interactions to target the resampling towards difficult-to-model dynamics, while~\citet{ziarko2026contrastive} utilizes different contexts for in-trajectory sampling. We also build on the successes of CRL when applied in simulated manipulation~\citep{wang20261000}, especially related to emergent exploration~\citep{liu2025single,bastankhah2026demystifying}, leveraging these properties which allow for limited manipulation performance in a sample-efficient manner relative to prior GCRL work~\citep{puterman1990markov,kaelbling1993learning,liu2022goal,andrychowicz2017hindsight,bai2021addressing}. Besides contrastive learning as applied to GCRL, \methodshort{} draws on ideas from general contrastive learning, especially related to leveraging privileged information such as class labels (wherein our case we leverage interactions and in-trajectory sampling)~\citep{feng2022adaptive,denize2023similarity,khosla2020supervised,hoffmann2022ranking}, and debiasing strategies which carefully select negative examples, which we adapt for interaction-based sampling~\citep{chuang2020debiased,huynh2022boosting,dwibedi2021little}. Our observations about the relationship between non-smooth dynamics and smooth representations learned by contrastive methods are in part inspired by~\citet{betser2026infonce}, which observes that the distribution of representations tends towards a thin-shell Gaussian. Finally, recent work has shown close connections between contrastive RL methods and the general class of methods used in unsupervised Reinforcement Learning~\citep{agarwal2025unified}, including learning to represent all rewards through successor representations~\citep{sikchi2026rlzero,zheng2026can,touati2021learning,agarwal2024proto,zheng2024contrastive}, mutual information skill learning~\citep{modirshanechi2026unifying,levy2025unsupervised}, representation learning~\citep{ferns2011bisimulation,zhang2020learning,rudolph2024learning} and world models~\citep{farebrother2026compositional}. However, the application of factorization and especially entity interactions to the methods based on contrastive learning, which we use in \methodshort{}, remains limited. 

\textbf{Manipulation Learning with Interactions}: Robot manipulation encompasses a wide range of tasks, from grasping and moving objects to striking, inserting or throwing~\citep{han2023survey}. Interactions~\citep{chuck2024automated} have been incorporated into learning for manipulation in a variety of ways, including controllability~\citep{seitzer2021causal}, data augmentation~\citep{pitis2020counterfactual,pitis2022mocoda}, hierarchical reinforcement learning~\citep {chuck2020hypothesis,chuck2023granger,wang2024skild}, skill learning~\citep{hu2024disentangled,rodriguez2025pixels,hosseini2026susd}, causal modeling~\citep{lee2023scale,biswas2024gaze} and exploration~\citep{wang2023elden}. This work introduces a novel formalization for the challenges of robot learning for manipulation as mode-changing dynamics, and \methodshort{} is a novel algorithm for incorporating interaction information into the training of contrastive RL. \methodshort{} shows significantly improved sample efficiency, and the first application to robot air hockey~\citep{chuck2024air}. Other work in manipulation that is more loosely leveraging interactions includes work centered around modeling dynamics, such as through graph networks~\citep{sieb2020graph,huang2023planning,lin2022efficient,zhang2025particle,huangrobocraft,kedia2024interact}, (which incorporate interactions through the message passing framework), human conditional intent~\citep{zhang2024adaptigraph} or human video~\citep{singh2025hand,zhu2026vision}, and affordances~\citep{khazatsky2021can,girgin2024multiobject}. Most closely related is the work investigating dynamics modeling~\citep{haramati2026hierarchical,wang2026factored} which learn sturctures that incorporate factors for policy learning, often also identifying the factors themselves from observations~\citep{daniel2026latent, qi2025ec}. \methodshort{} is applied to CRL, which learns structure for GCRL policies, but these model-learning algorithms are likely to be effective in concert with CRL methods, where the dynamics model and policies can work together.